\teachcover{Explaining\\Fibonacci periods}{1}{
The field containing the roots controls the powers
\[
\pi(11)=10,\qquad\pi(7)=16.
\]
}{
The first three lectures supplied fields, coordinates, cyclic
multiplicative groups, and counts of GL and SL. Today returns to
Fibonacci with a full 50 minutes for the order argument and Frobenius.
The prime binomial identity and Fermat's theorem are the remaining
prerequisites. Keep the calculation break before the modulo-11 example.
}

\teachslide{The structures we now understand}{2}{
For a finite field $K$ with $q$ elements, $K^\times$ is cyclic of order $q-1$.
\[
|\SL_2(\F_p)|=p(p^2-1).
\]
For Fibonacci, with $A=\mat1110$,
\[
\det A=-1,\qquad A^2=A+I.
\]
\key{For odd $p$, apply Lagrange to $A^2\in\SL_2(\F_p)$.}
}{
Recall that a primitive element has order $q-1$; not every element
is primitive. Matrix groups need not be cyclic. The determinant
of $A$ is -1, so its square rather than $A$ itself lies in SL.
Our field-group theorem is about multiplication of nonzero scalars.
}

\teachslide{The Fibonacci period is the matrix order}{4}{
For $t\ge1$,
\[
A^t=\begin{pmatrix}F_{t+1}&F_t\\F_t&F_{t-1}\end{pmatrix}.
\]
Since $F_{t-1}=F_{t+1}-F_t$,
\[
A^t=I\iff(F_t,F_{t+1})=(0,1)\pmod p.
\]
\[
\key{\pi(p)=\ord(A).}
\]
}{
The displayed power formula follows by induction from the recurrence,
and is the students' existing matrix approach. The equivalence needs
both state coordinates, not merely $F_t=0$. If a student asks why an
orbit period should equal the whole matrix order, this identity is
the specific Fibonacci justification; that equality need not hold for
an arbitrary matrix and arbitrary starting vector.
}

\teachslide{A first bound from $\SL_2$}{2}{
Let $p$ be odd. If $A^t=I$, determinants give $(-1)^t=1$.

Thus $\pi(p)$ is even, and
\[
\ord(A^2)=\frac{\pi(p)}2.
\]
Lagrange gives
\[
\key{\pi(p)\mid2p(p^2-1).}
\]
Can the roots of $X^2-X-1$ give a sharper answer?
}{
For an element of even order $t$, the least positive $r$ with $A^{2r}=I$
is $t/2$, because $t$ must divide $2r$. This proves the order identity.
The broad group-size divisibility uses no splitting information.
Its purpose is to motivate exploiting the special matrix, rather
than presenting the group size as an exact Fibonacci period.
}

\teachslide{Choose a field containing the roots}{3}{
Assume $p\ne2,5$.
\begin{itemize}
\item If $5$ is square in $\F_p$, use $K=\F_p$.
\item Otherwise, use $K=\{a+bs:s^2=5\}$ over $\F_p$.
\end{itemize}
\[
\alpha=\frac{1+s}{2},\qquad\beta=\frac{1-s}{2}.
\]
\[
\alpha+\beta=1,\quad\alpha\beta=-1,\quad\alpha\ne\beta.
\]
}{
In the split case, $s$ is an existing square root, not a new symbol.
In the nonsplit case, the proof from Lecture 2 generalizes directly
from $7$ to any odd $p$: $a^2-5b^2$ cannot vanish for a nonzero pair.
The prime $5$ is excluded because the roots coincide, and $2$ because
division by $2$ is unavailable. Both roots are nonzero since their
product is $-1$.
}

\teachslide{Diagonalize in the field where it makes sense}{4}{
For either root $r$, since $r^2=r+1$,
\[
A\begin{pmatrix}r\\1\end{pmatrix}
=r\begin{pmatrix}r\\1\end{pmatrix}.
\]
\[
P=\begin{pmatrix}\alpha&\beta\\1&1\end{pmatrix},
\qquad \det P=\alpha-\beta\ne0.
\]
\[
P^{-1}AP=\operatorname{diag}(\alpha,\beta),\qquad
\key{A^t=I\iff\alpha^t=\beta^t=1.}
\]
}{
Compute the product $(r+1,r)^{\mathsf T}$ and use the root equation.
This explicitly proves the eigenvectors and invertibility of the
eigenvector matrix, so no general diagonalization theorem is needed.
Changing the field lets us find $P$; it does not change whether a
power of the original matrix equals $I$. Its entries still lie in
the embedded prime field. This also gives the lcm of the two root orders.
}

\teachslide{A period is a simultaneous order}{2}{
We have shown
\[
A^t=I\iff\alpha^t=\beta^t=1.
\]
If $\ord(\alpha)=r$ and $\ord(\beta)=s$, the return times are
the common multiples of $r$ and $s$.
\[
\key{\pi(p)=\operatorname{lcm}(r,s).}
\]
}{
An element's return times are exactly the multiples of its order,
by division with remainder. We need both eigenvalues to return to 1.
This lcm statement is about simultaneous return, not the order of
the product of arbitrary elements. The coprime-product lemma from
Lecture 3 had an additional coprimality hypothesis.
}

\teachslide{Case one: the roots are in $\F_p$}{3}{
If $5$ is a square in $\F_p$, both roots lie in $\F_p^\times$.

Fermat's theorem gives
\[
\alpha^{p-1}=\beta^{p-1}=1.
\]
Therefore
\[
\key{A^{p-1}=I,\qquad\pi(p)\mid p-1.}
\]
}{
Remind students that the roots are nonzero, so Fermat applies.
The conclusion is divisibility, not equality. The new bound is
much sharper than the broad SL group-size bound because we have
used the eigenvalues of this particular matrix.
}

\teachslide{Your turn: roots modulo eleven}{3}{
In $\F_{11}$, use $s=4$ and $2^{-1}=6$.
\begin{enumerate}
\item Compute $\alpha=(1+s)/2$ and $\beta=(1-s)/2$.
\item Compute $\alpha^2$ and $\alpha^5$ by repeated squaring.
\item Which proper divisors of $10$ could still be the period?
\end{enumerate}
\spacer
\key{We already know that the period divides $10$.}
}{
Give students a minute to calculate before discussing. Expected
roots are 8 and 4; $8^2=9$, $8^4=4$, $8^5=-1$. These rule out
both 2 and 5 as return times, and therefore all proper divisors of
10. The next page reveals the answers and exact period.
}

\teachslide{Modulo eleven: an exact period}{2}{
Choose $s=4$ in $\F_{11}$. Then
\[
\alpha=8,\qquad\beta=4,\qquad\pi(11)\mid10.
\]
But
\[
8^5=-1,\qquad8^2=9\ne1\pmod{11}.
\]
Every proper divisor of $10$ divides $5$ or $2$.
\[
\key{\pi(11)=10.}
\]
}{
Compute $2^{-1}=6$ modulo $11$, giving roots $5\cdot6=8$ and
$(-3)\cdot6=4$. To check the powers, $8^2=9$, $8^4=4$, and
$8^5=32=-1$. If $A^5$ or $A^2$ were $I$, the corresponding power
of $8$ would be $1$. We have ruled out every proper divisor of the
known multiple $10$. The order of the other root need not be computed.
}

\teachslide{A Fibonacci root need not be primitive}{3}{
Modulo $11$, $\alpha=8$ has order $10$ and generates $\F_{11}^\times$.

Modulo $7$, the roots live in $\F_{49}$, where
\[
|\F_{49}^\times|=48.
\]
But $\alpha=4+4s$ satisfies $\alpha^8=-1$ and $\alpha^{16}=1$.
\[
\key{\ord(\alpha)=16,\quad\text{so it is not a multiplicative generator}.}
\]
}{
This comparison connects cyclicity directly to the order problem.
In the field from Lecture 2, $\alpha^2=5+4s$, $\alpha^4=5s$,
and $\alpha^8=(5s)^2=125=-1$. Thus its exact order is 16.
Nevertheless $s=2\alpha-1$, so $\alpha$ does generate this field
over $\F_7$. Generating the extension field and generating its
multiplicative group are different requirements. Frobenius will
explain the order bound without needing these special calculations.
}

\teachslide{A new symmetry: Frobenius}{4}{
In characteristic $p$,
\[
\key{(u+v)^p=u^p+v^p.}
\]
The intermediate binomial coefficients are divisible by $p$.

Also $(uv)^p=u^pv^p$ and $1^p=1$.

On a finite field, $u\mapsto u^p$ is injective, hence bijective:
\[
u^p=v^p\quad\Longrightarrow\quad(u-v)^p=0\quad\Longrightarrow\quad u=v.
\]
}{
Recall why $p$ divides $\binom pj$ for $0<j<p$: the numerator
contains a factor $p$, while $j!$ is prime to $p$. The power map
preserves the field operations and is a bijection, so it is a field
automorphism, called Frobenius. A nonzero field element cannot have
a zero power. Finiteness is needed for the injective-to-surjective
step; no infinite-field claim is required.
}

\teachslide{Which elements does Frobenius fix?}{3}{
Every $a\in\F_p$ satisfies $a^p=a$.

The fixed elements are roots of
\[
X^p-X.
\]
A nonzero polynomial of degree $p$ has at most $p$ roots.

The prime field already supplies $p$ roots.
\[
\key{u^p=u\iff u\in\F_p.}
\]
}{
For nonzero prime-field elements use Fermat; zero is immediate.
Recall the polynomial root bound proved in Lecture 3. The polynomial is nonzero
since its leading coefficient is $1$, even in characteristic $p$.
This argument identifies the fixed elements in our extension without
assuming the desired conjugation rule.
}

\teachslide{Case two: Frobenius exchanges the roots}{3}{
Suppose $5$ is nonsquare in $\F_p$, with $s^2=5$ in $K$.
\[
(s^p)^2=5^p=5.
\]
Thus $s^p$ is $s$ or $-s$.

It cannot be $s$, because $s\notin\F_p$.
\[
\key{s^p=-s,\qquad\alpha^p=\beta,\qquad\beta^p=\alpha.}
\]
}{
The factorization $X^2-5=(X-s)(X+s)$ in $K$ and absence of zero
divisors give exactly those two square roots. They are distinct since
$p$ is odd and $s\ne0$. If $s^p=s$, the preceding fixed-element result
would put $s$ in the prime field, contradicting nonsquareness.
Frobenius fixes $1$ and $2^{-1}$, so applying it to the formulas for
$\alpha,\beta$ interchanges them.
}

\teachslide{The minus sign tells us the return time}{3}{
\[
\alpha^{p+1}=\alpha\beta=-1,
\qquad\beta^{p+1}=-1.
\]
Diagonalization gives
\[
\key{A^{p+1}=-I,\qquad A^{2(p+1)}=I.}
\]
Therefore
\[
\pi(p)\mid2(p+1),\qquad\pi(p)\nmid p+1.
\]
At time $p+1$, $(F_{p+1},F_{p+2})=(0,-1)$.
}{
Conjugating the scalar matrix $-I$ leaves it unchanged. Since the
characteristic is odd, $-I\ne I$. Applying the matrix to $(1,0)$
gives $(F_{p+2},F_{p+1})=(-1,0)$, hence the displayed state in the
usual order. This connects the minus sign to the earlier warning
that a Fibonacci zero alone is not a restart.
}

\teachslide{Modulo seven: the other exact period}{3}{
Since $5$ is nonsquare modulo $7$,
\[
A^8=-I,\qquad A^{16}=I.
\]
Thus $\pi(7)\mid16$ but $\pi(7)\nmid8$.

Every proper divisor of $16$ divides $8$.
\[
\key{\pi(7)=16.}
\]
}{
This exact-order argument avoids a long power calculation in the
extension. Ask students why $A^8\ne I$ rules out every proper divisor
of $16$. For a direct Frobenius check in the field constructed last
time, $s^7=s(s^2)^3=125s=-s$. The general proof explains why this
calculation is not a coincidence.
}

\teachslide{What the algebra explained}{2}{
For primes $p\ne2,5$,
\[
\begin{array}{c|c}
5\text{ is square in }\F_p&\pi(p)\mid p-1\\\hline
5\text{ is nonsquare in }\F_p&\pi(p)\mid2(p+1)
\end{array}
\]
\begin{itemize}
\item Group size supplies a broad order bound.
\item The field containing the roots supplies a sharper bound.
\item Further power checks identify the least period.
\end{itemize}
}{
Recap the distinction between existence of a period, a known multiple,
and the least period. The matrix construction proves existence for
all moduli; group structure and the root field give increasingly
specific information for prime moduli. Prime $29$ is an optional
example where the bound is not attained.
}

\teachslide{Exit question}{3}{
Suppose $p$ is odd and $5$ is nonsquare modulo $p$.

A student argues: ``$F_{p+1}=0$, so the period is $p+1$.''

What is missing, and what does the matrix actually tell us?
\spacer
\key{Use both entries of the Fibonacci state.}
}{
Expected answer: a restart requires $(F_{p+1},F_{p+2})=(0,1)$,
but the state is $(0,-1)$ because $A^{p+1}=-I$. The period divides
$2(p+1)$ and does not divide $p+1$; it need not equal $2(p+1)$.
End the main lecture here at minute 50. Use optional slides only
if time remains or as material for a subsequent meeting.
}

\backups
\backupslide{Why the prime five is different}{
Modulo $5$, $X^2-X-1=(X-3)^2$.
\[
A=3I+N,\qquad N=\mat3112,\qquad N^2=0.
\]
For $t\ge1$,
\[
A^t=3^tI+t3^{t-1}N.
\]
$A^t=I$ requires $5\mid t$ and $4\mid t$.
\[
\key{\pi(5)=20.}
\]
}{
Compute $A-3I$ modulo $5$ to get the displayed nonzero nilpotent matrix.
The off-diagonal entry in $A^t$ forces $t=0$ in the field, hence
$5\mid t$, because $3$ is nonzero. Then $A^t=3^tI$ and $3$ has
order $4$ modulo $5$: $3^2=-1$ and $3^4=1$. These two divisibilities
are also sufficient, giving least time $20$. The repeated root
produces a nilpotent term instead of two independent eigenvectors.
}

\backupslide{The prime two}{
Modulo $2$, the successive Fibonacci states are
\[
(0,1)\ \longmapsto\ (1,1)\ \longmapsto\ (1,0)
\ \longmapsto\ (0,1).
\]
\[
\key{\pi(2)=3.}
\]
Here $-1=1$, and division by $2$ is unavailable.
}{
The short state cycle gives the exact answer immediately. This also
explains why the odd-characteristic parity argument and the root
formula with denominator $2$ were explicitly restricted.
}

\backupslide{A bound need not be the least period}{
Modulo $29$, $11^2=5$, so $\pi(29)\mid28$.

The roots are $\alpha=6$ and $\beta=24$.
\[
6^7=-1,\quad6^2=7\ne1;
\qquad24^7=1,\quad24\ne1.
\]
Their orders are $14$ and $7$.
\[
\key{\pi(29)=\operatorname{lcm}(14,7)=14.}
\]
}{
Check $6^7=-1$ by repeated squaring: $6^2=7$, $6^4=20$,
$6^7=20\cdot7\cdot6=28$. Thus $6$ has order $14$, since every
proper divisor of $14$ divides $7$ or $2$. The other root is
$-6^{-1}$, so $24^7=(-1)^7(6^7)^{-1}=1$; since $7$ is prime
and $24\ne1$, its order is $7$. This is a concrete proper-divisor
example, not just a caution about the word divides.
}

\backupslide{An exact-period procedure}{
Start with a proven multiple $B$ of $\pi(p)$.
\begin{enumerate}
\item Set $t=B$ and factor $t$.
\item For each prime divisor $q$ of $t$, test $A^{t/q}=I$.
\item If the test succeeds, replace $t$ by $t/q$ and repeat.
\end{enumerate}
When no prime factor can be removed, $t=\pi(p)$.
}{
Matrix repeated squaring performs the tests efficiently. At every
stage the true order divides $t$. If $t$ were still larger than
the order, some prime factor of $t/\pi(p)$ could be removed while
leaving a multiple of the order. Thus the stopping rule is exact.
For $29$, start at $28$ and remove one factor $2$ to reach $14$.
}
